% Хэрэглээний тохиолдол 1: Шошго үүсгэх
\begin{longtable}{|l|p{10.5cm}|}
    \caption{Хэрэглээний тохиолдол 1: Шошго үүсгэх} \label{tab:uc_create_category} \\
    \hline
    \textbf{ID / Нэр} & UC-01 / Шинэ шошго үүсгэх \\ \hline
    \textbf{Товч тайлбар} & Хэрэглэгч зурвас шүүх шинэ дүрэм (шошго) үүсгэж, нэр, өнгө болон шүүлтүүрүүдийг тохируулна. \\ \hline
    \textbf{Үйлдэгч} & Хэрэглэгч \\ \hline
    \textbf{Өмнөх нөхцөл} & Хэрэглэгч ``Шошго удирдах'' дэлгэцийг нээсэн байна. \\ \hline
    \textbf{Ажлын урсгал} &
    1. Хэрэглэгч ``+'' товчийг дарж үүсгэх цонхыг нээнэ. \newline
    2. Шошгоны нэрийг оруулж, жагсаалтаас өнгө сонгоно. \newline
    3. Шүүлтүүрийн төрлийг сонгоно: \newline
    \hspace{0.5cm} 3.1. Энгийн үг оруулах бол ``Энгийн үг нэмэх'' талбарт бичнэ. \newline
    \hspace{0.5cm} 3.2. Тогтмол илэрхийлэл (Regex) ашиглах бол дүрс дээр дарж паттерн оруулна. \newline
    4. Оруулсан өгөгдлийг шалгаж ``OK'' товчийг дарна. \\ \hline
    \textbf{Альтернатив} & Хэрэглэгч ``Cancel'' дарвал өөрчлөлт хадгалагдахгүй. Буруу тогтмол илэрхийллийн паттерн оруулбал систем алдаа заана. \\ \hline
    \textbf{Дараах нөхцөл} & Шинэ шошго баазад хадгалагдаж, зурвасууд шууд ангилагдана. \\ \hline
\end{longtable}

\vspace{1em}

% Хэрэглээний тохиолдол 2: Хавтас үүсгэх
\begin{longtable}{|l|p{10.5cm}|}
    \caption{Хэрэглээний тохиолдол 2: Хадгалсан харагдац (Хавтас) үүсгэх} \label{tab:uc_create_folder} \\
    \hline
    \textbf{ID / Нэр} & UC-02 / Хавтас үүсгэх \\ \hline
    \textbf{Товч тайлбар} & Хэрэглэгч өөрийн хэрэгцээнд тохирсон зурвасын харагдац (Хавтас) үүсгэж, доод навигацийн цэсэнд нэмнэ. \\ \hline
    \textbf{Үйлдэгч} & Хэрэглэгч \\ \hline
    \textbf{Ажлын урсгал} &
    1. Баруун дээд цэснээс ``Хавтас'' -> ``Харагдац үүсгэх'' сонгоно. \newline
    2. Хавтасны нэрийг бичиж, харагдах байрлал (index) болон өнгийг тодорхойлно. \newline
    3. ``OK'' товчийг дарж баталгаажуулна. \\ \hline
    \textbf{Дараах нөхцөл} & Шинэ хавтас доод навигацийн цэсэнд нэмэгдэж, хэрэглэгч тухайн хавтас руу автоматаар шилжинэ. \\ \hline
\end{longtable}

\vspace{1em}

% Хэрэглээний тохиолдол 3: Шошгоор шүүх
\begin{longtable}{|l|p{10.5cm}|}
    \caption{Хэрэглээний тохиолдол 3: Шошгоор шүүх} \label{tab:uc_filter_category} \\
    \hline
    \textbf{ID / Нэр} & UC-03 / Харилцан яриаг шошгоор шүүх \\ \hline
    \textbf{Товч тайлбар} & Хэрэглэгч үүсгэсэн хавтас доторх харилцан яриануудыг нэг болон хэд хэдэн шошгоор (Tag) шүүж харна. \\ \hline
    \textbf{Ажлын урсгал} &
    1. Хайх талбарын хажууд байрлах шүүлтүүрийн товчийг дарна. \newline
    2. Гарч ирэх жагсаалтаас нэг болон түүнээс дээш шошгыг сонгоно. \newline
    3. ``OK'' товчийг дарахад систем сонгосон шошгонд таарсан зурвасуудыг харуулна. \\ \hline
    \textbf{Дараах нөхцөл} & Жагсаалт шинэчлэгдэж, шүүлтийн төлөв хадгалагдана. \\ \hline
\end{longtable}

\vspace{1em}

% Хэрэглээний тохиолдол 4: Шудрах үйлдэл
\begin{longtable}{|l|p{10.5cm}|}
    \caption{Хэрэглээний тохиолдол 4: Шудрах үйлдлийн тохиргоо} \label{tab:uc_swipe_actions} \\
    \hline
    \textbf{ID / Нэр} & UC-04 / Шудрах үйлдлийн тохиргоо \\ \hline
    \textbf{Товч тайлбар} & Харилцан ярианы жагсаалтад зүүн эсвэл баруун тийш шудрахад хийгдэх үйлдлийг хэрэглэгч өөрөө сонгоно. \\ \hline
    \textbf{Ажлын урсгал} &
    1. Програмын ерөнхий тохиргооноос ``Шудрах үйлдэл'' хэсгийг нээнэ. \newline
    2. ``Эхнээс нь шудрах'' эсвэл ``Төгсгөлөөс нь шудрах'' сонголтыг дарна. \newline
    3. Боломжит үйлдлүүдээс (Устгах, Архивлах, Уншсан болгох г.м) нэгийг сонгоно. \\ \hline
    \textbf{Дараах нөхцөл} & Сонгосон үйлдэл тэр даруй хадгалагдаж, жагсаалт дээр ажиллаж эхэлнэ. \\ \hline
\end{longtable}

\vspace{1em}

% Хэрэглээний тохиолдол 5: Дэлгэцийн горим
\begin{longtable}{|l|p{10.5cm}|}
    \caption{Хэрэглээний тохиолдол 5: Дэлгэцийн харагдацын горим} \label{tab:uc_screen_view} \\
    \hline
    \textbf{ID / Нэр} & UC-05 / Screen View Mode сонгох \\ \hline
    \textbf{Товч тайлбар} & Нугалдаг төхөөрөмж эсвэл таблет дээр жагсаалт болон дэлгэрэнгүй хэсгийг хэрхэн харуулахыг тохируулна. \\ \hline
    \textbf{Ажлын урсгал} &
    1. Тохиргоо цэсний ``View mode'' хэсэгт орно. \newline
    2. ``Auto'', ``Single'', эсвэл ``Separated'' горимуудаас нэгийг сонгоно. \newline
    3. ``Auto'' горимд систем дэлгэцийн төлөвийг (Дэлгэц хаалттай/нээлттэй) мэдэрч автоматаар шилжинэ. \\ \hline
    \textbf{Дараах нөхцөл} & Програмын үндсэн дэлгэцийн бүтэц (layout) сонгосон горимын дагуу шинэчлэгдэнэ. \\ \hline
\end{longtable}

\vspace{1em}

% Хэрэглээний тохиолдол 6: Хавтас засварлах
\begin{longtable}{|l|p{10.5cm}|}
    \caption{Хэрэглээний тохиолдол 6: Хавтас засварлах ба устгах} \label{tab:uc_edit_folder} \\
    \hline
    \textbf{ID / Нэр} & UC-06 / Хавтас засах эсвэл устгах \\ \hline
    \textbf{Товч тайлбар} & Өмнө үүсгэсэн хавтасны мэдээллийг шинэчлэх эсвэл системээс бүрмөсөн устгана. \\ \hline
    \textbf{Ажлын урсгал} &
    1. Доод навигацийн цэс дээрх хавтасны дүрс дээр удаан дарна. \newline
    2. Мэдээллийг (нэр, өнгө, байрлал) засаад ``OK'' дарах эсвэл ``Delete'' товчийг сонгоно. \newline
    3. Устгах үйлдэл дээр систем дахин баталгаажуулалт асууна. \\ \hline
    \textbf{Дараах нөхцөл} & Өөрчлөлт хадгалагдах эсвэл хавтас навигацийн цэснээс алга болж хэрэглэгч үндсэн харагдац руу шилжинэ. \\ \hline
\end{longtable}
